\documentstyle[% 


righttag% 


,11pt% 


,amssymb% 


,russian%               remove in Eng. 


,izvru%                 Set izven% in Eng. 


]{amsart} 


 


% 


 


\newcommand{\const}{\operatorname{const}} 


\newcommand{\tg}{\operatorname{tg}} 


% \newcommand{\arcsin}{\operatorname{arcsin}} 


\newcommand{\Grad}{\operatorname{grad}} 


\newcommand{\tts}[1]{} 


 


\theoremstyle{plain} 


\theorembodyfont{\it} 


\newtheorem{thmnonum}{\hskip\parindent Теорема} 


\renewcommand{\thethmnonum}{} 


 


 


%\hoffset=-15mm%                  For Eng. 


\setcounter{page}{29} 


\god{2002} 


\nomer{4~(479)} %                        Remove in Eng. 


%\nomer{Vol.\,46, No.\,4}%               Set in Eng. 


%\str{\pageref{begin}--\pageref{end}}%  Set in Eng. 


\udk{517.544} 


 


\author{\it И.Т.\,Денисюк} 


% NB:   ^^ remove in Eng. 


\title{Задача сопряжения гармонических функций\\ 


в трехмерных областях с негладкими границами} 


 


\date{} 


%\pagestyle{myheadings}%         Set in Eng. 


\pagestyle{plain} %              Remove in Eng. 


\begin{document} 


%\thispagestyle{plain}%          Set in Eng. 


\maketitle 


\label{begin} 


 


 


Задачи сопряжения аналитических функций в двумерных областях с  


кусочно-гладкими границами изучались в работах [1], [2]. В данной 


статье  


рассматривается задача построения гармонических функций в областях  


пространства $R^3$ при условии сопряжения их на поверхностях, содержащих  


негладкости типа особых линий (множеств угловых точек) и конических 


точек. 


\medskip 


 


{\bf 1.} {\it Постановка задачи}. Пусть в $R^3$ содержатся односвязные 


конечные области 


$V_i$, $i=\ov{1,k_0}$, $V_i\cap V_l=\emptyset$, ограниченные замкнутыми 


поверхностями 


$S_i$, содержащими гладкие замкнутые непересекающиеся особые линии 


(множества угловых точек) $L_{i j}$, $j=\ov{1,p_i}$, и конические точки 


$O_{i k}\notin L_{i j}$, $k=\ov{1,q_i}$. 


 


Построим функции $F_i(x,y,z)$, гармонические в соответствующих областях 


$V_i$, и функцию $F_0(x,y,z)$, гармоническую в 


$V_0=R^3\setminus\bigcup\limits^{k_0}_{i=1}V_i$, при таких условиях 


сопряжения на поверхностях~$S_i$: 


 


в точках гладкости $N_0\in S_i$, $N_0\notin L_{i j}$, $N_0\ne O_{i k}$, 


\begin{gather} 


F^+_i(N_0)-F^-_0(N_0)=0,\\ 


\lambda_i\partial F^+_i(N_0)/\partial n- 


\lambda_0\partial F^-_0(N_0)/\partial n=0; 


\end{gather} 


 


в точках особых линий $L_{i j}$ или в конических точках $O_{i k}$ 


\begin{equation} 


\lim[F^+_i(N_0)-F^-_0(N_0)]=0,\quad 


\lim[\lambda_i\partial F^+_i(N_0)/\partial n-\lambda_0 


\partial F^-_0(N_0)/\partial n]=0, 


\end{equation} 


где пределы берутся при $N_0\to M_0\in L_{i j}$ в плоскости, нормальной 


к особой 


линии, или при $N_0\to O_{i k}$, постоянные $\lambda_0,\lambda_i\in R$, 


$\bold n$ --- внешняя нормаль к $S_i$. 


Верхний индекс $\pm$ означает граничное значение гармонической функции 


при  


подходе к поверхности $S_i$ со стороны области $V_i$ (знак~$+$) или 


$V_0$ 


(знак~$-$). Считаем, что вблизи своей вершины коническая поверхность 


имеет  


прямолинейную образующую и криволинейную гладкую замкнутую  


направляющую. 


 


К такой задаче приводятся некоторые проблемы механики сплошной 


неоднородной  


среды [3], [4]. 


Гармоничность функций $F_i(M)$ в точках особых линий $L_{i j}$ или 


в конических точках $O_{i k}$ понимается как выполнение равенства 


\begin{equation} 


\lim\iiint\limits_{V^*}\nabla^2F(M)d v=0,\quad M^*_0\in V^*, 


\end{equation} 


где пределы берутся при $V^*\to M^*_0$, т.\,е. при стягивании области 


$V^*$ в точку  


$M^*_0$, что физически обусловлено необходимостью выполнения законов  


сохранения энергии в особых точках среды [3]--[5], при этом 


$M=M(x,y,z)$, 


$M^*_0\in L_{i j}$ или $M^*_0=O_{i k}$. 


\medskip 


 


{\bf 2}. {\it Асимптотика гармонических функций}. Для установления 


характера 


гармонических функций изучим их поведение в локальных областях особых 


точек. 


 


\begin{lem} 


Если в точках особой линии $L$ величина 


\begin{equation} 


m=\pi/(2\pi-\varpi), 


\end{equation} 


где $\varpi$ --- угол раствора поверхностей в точках особой линии в 


плоскости, 


нормальной к ней, принадлежит интервалу $(0,1)$, то гармоническая 


функция 


и ее нормальная производная имеют асимптотическое представление 


\begin{equation} 


\begin{array}{c} 


F_i(x,y,z)=\rho^m(C_{i1}\cos m\theta+C_{i2}\sin m\theta)+ 


O(\rho^{m+1}),\\[2mm] 


\partial F_i(x,y,z)/\partial n=\rho^{m-1} 


(C_{i2}m\cos m\theta-C_{i1}m\sin m\theta) 


+O(\rho^m),\ \ i=0,1, 


\end{array} 


\end{equation} 


где $\rho$, $\theta$, $s$ --- локальные координаты, $C_{i q}=C_{i q}(s)$, 


$q=1,2$, --- величины, зависящие от $s$. 


\end{lem} 


 


\begin{pf} 


Пусть особая линия $L$ является пересечением гладких  


поверхностей $S_1$ и $S_2$, $L=S_1\cap S_2$. Натянем на $L$ гладкую 


поверхность $S_0$, на 


которой введем ортогональные координаты $u$, $v$ так, что при 


$v=\const$ 


описывается особая линия $L$, и в результате параметризации полагаем 


$u=s$, где  


$s$ --- длина дуги $L$, отсчитываемая от некоторой начальной точки 


([6], с.\,58). 


Пусть $\bold n\bold n_1\bold n_2$ --- сопровождающий трехгранник 


поверхности $S_0$ в точке $M_0$ кривой, 


где $\bold n_1$, $\bold n_2$ --- орты, лежащие в плоскости, касательной 


к $S_0$ в точке 


$M_0$, $\bold n_1$ --- касательный вектор к кривой $L$, $\bold 


n\bot\bold n_1$, $\bold n\bot\bold n_2$. Точки $M$ плоскости 


векторов $\bold n$, $\bold n_2$ определяем полярными координатами 


$(\rho,\theta)$, $\theta$ --- угол между 


векторами $\bold M_0\bold M$ и $\bold n_2$. 


 


Представление 


\begin{equation} 


\bold r=\bold r_0+\rho(\cos\theta\bold n_2(s)+\sin 


\theta\bold n(s)), 


\end{equation} 


где $\bold r_0$, $\bold r$ --- радиус-векторы точек $M_0$, $M$, 


содержит 


криволинейные ортогональные координаты $\rho$, $\theta$, $s$ точки $M$ 


локальной области  


пространственной кривой $L$ с коэффициентами Ламе $h_1=1$, $h_2=\rho$, 


$h_3=1+\rho H_0$, $H_0=-(p\cos\theta+r_1\sin\theta)$, $p=-|\bold 


r_s|_v/|\bold r_v|$, $r_1=(\bold n,\bold r_{s s})/|\bold r_s|$. 


 


Уравнение Лапласа $\Delta F(x,y,z)=0$ 


допускает преобразование подобия 


\begin{equation} 


F=A_1F(B_1x,B_1y,B_1z), 


\end{equation} 


где $A_1=A_1(B_1)$, $A_1,B_1\in R$. 


Дифференцируя соотношение (8) по $B_1$, получаем уравнение 


$$ 


F(x,y,z)d A_1/d B_1+A_1(\Grad F(x,y,z),\bold r)=0, 


$$ 


в решение которого 


\begin{equation} 


F(x,y,z)=F^*_0(x,y,z) 


\end{equation} 


входит однородная функция $F^*_0(x,y,z)$\; $m$-го порядка переменных 


$x$, $y$, $z$. Переходя 


в (9) к криволинейным координатам (7), устанавливаем, что гармоническая  


функция принадлежит степенному классу относительно переменной $\rho$ 


\begin{equation} 


F_i=\rho^m A_i(\theta,s),\quad i=0,1. 


\end{equation} 


Подставляя представление (10) в уравнение (4), записанное в координатах  


$\rho$, $\theta$, $s$, находим определяющее уравнение для функции 


$A_i(\theta,s)$ 


$$ 


\partial^2A_i/\partial\theta^2+m^2A_i=0,\quad m>0, 


$$ 


решение которого и дает асимптотическое представление (6). Удовлетворяя  


граничным условиям (3) с помощью (10), получаем однородную систему  


линейных алгебраических уравнений относительно величин $C_{i q}$. Она 


разрешима, если ее определитель равен нулю, что и дает соотношение (5). 


\end{pf} 


 


Пусть точка $O$ является вершиной конической поверхности с  


прямолинейными образующими и гладкой замкнутой направляющей $L_0$, 


очерченной ортом $\bold r_0(s_1)$ образующей с началом в точке $O$. 


Здесь $s_1$ --- длина дуги  


направляющей, отсчитываемая от некоторой начальной точки. Введем 


локальные  


криволинейные координаты $\rho_1$, $\theta_1$, $s_1$ согласно равенству 


\begin{equation} 


\bold r_1=\rho_1(\bold r_0(s_1)\sin\theta_1+ 


\bold n_3(s_1)\cos\theta_1), 


\end{equation} 


где $\bold r_1$ --- радиус-вектор точки $M(x,y,z)$, $\bold 


n_3(s_1)=[\bold r_0(s_1),d\bold r_0(s_1)/d s_1]$ --- нормаль к  


образующей $\bold r_0(s_1)$ конической поверхности в точке $O$, 


$\rho=|\bold r_1|$, $\theta_1$ --- угол,  


образованный вектором $\bold r_1$ с нормалью $\bold n_3(s_1)$. Такие 


криволинейные координаты 


ортогональны с коэффициентами Ламе $h_1=1$, $h_2=\rho_1$, 


$h_3=\rho_1H_{01}$, $H_{01}=\sin\theta_1+d_1\cos\theta_1$, $d_1=|[\bold 


r_0(s_1),d^2\bold r_0(s_1)/d s^2_1]|$, 


а при $\theta_1=\pi/2$ уравнение (11) описывает коническую поверхность. 


Аналогично  


доказательству леммы 1 устанавливаем, что гармоническая функция в  


окрестности конической точки принадлежит степенному классу 


\begin{equation} 


F(x,y,z)=\rho^{m_0}_1D(\theta_1,s_1). 


\end{equation} 


Удовлетворяя уравнению (4), записанному в координатах $\rho_1$, 


$\theta_1$, $s_1$, с помощью  


(12) находим уравнение 


\begin{equation} 


Q_{m_0}[D(\theta_1,s_1)]=0, 


\end{equation} 


где оператор имеет вид 


$$ 


Q_{m_0}[\cdot]=H^{-1}_{01}(\partial(H_{01}\partial/ 


\partial\theta_1)/\partial\theta_1+\partial(H^{-1}_{01}\partial/ 


\partial s_1)/\partial s_1)+m_0(m_0+1). 


$$ 


Величины $H_{01}$, $H^{-1}_{01}$ являются дифференцируемыми, 


периодическими по $s_{01}$ 


функциями ($s_{01}$ --- длина направляющей $L_0$) и в таком классе 


строим функцию ([7], с.\,507) 


\begin{equation} 


D(\theta_1,s_1)=\sum^K_{n=0}(D_{n1}(\theta_1)\cos n l s_1+ 


D_{n2}(\theta_1)\sin n l s_1)+O(1/K), 


\end{equation} 


где $l=2\pi/s_{01}$. 


 


Подставляя отрезки рядов Фурье для $H_0$, $H^{-1}_0$, $D(\theta_1,s_1)$ 


типа (14) в уравнение 


(13), группируя и приравнивая выражения при гармониках $n=\ov{0,K}$ 


нулю,  


получим систему $2K+1$ обыкновенных уравнений второго порядка  


относительно неизвестных функций $D_{n q}(\theta_1)$, $q=1,2$. Такая 


система совместима 


([8], сс.\,12, 135) и имеет решение 


\begin{equation} 


D_{n q}(\theta_1)=\sum^{2(2K+1)}_{k=1}C_{n q k}\Phi_k(\theta_1), 


\end{equation} 


где $\Phi_k(\theta_1)$ --- ограниченные и определенные при 


$\theta_1\in[-\pi,\pi)$ функции  


фундаментального множества решений системы, $C_{n q k}$ --- 


произвольные 


постоянные. 


 


Из (12), (15) следует представление гармонической функции в локальной  


области 


\begin{equation} 


F(\rho_1,\theta_1,s_1)=\rho^{m_0}_1\bigg(\sum^K_{n=0} 


(D_{n1}(\theta_1)\cos n l s_1+D_{n2}(\theta_1)\sin n l s_1)+ 


O(1/K)\bigg)+O(\rho^{m_0+1}_1). 


\end{equation} 


Если $m_0\in(0,1)$, то нормальная производная принимает сингулярности  


степенного характера в конической точке 


\begin{equation} 


\partial F/\partial n=\rho^{m_0-1}_1\bigg(\sum^K_{n=0} 


(d D_{n1}(\theta_1)/d\theta_1\cos n l s_1+ 


d D_{n2}(\theta_1)/d\theta_1\sin n l s_1)+O(1/K)\bigg)+ 


O(\rho^{m_0}_1). 


\end{equation} 


 


Пусть $H_{01}$ в коэффициентах Ламе соотношения (11) не зависит от  


переменной $s_1$, что отвечает, например, круговому конусу, тогда  


$H_{01}=(1+d^2)^{1/2}\sin\theta_{11}$, $\theta_{11}=\theta_1+\beta$, 


$\tg\beta=d$. При этом функция $D(\theta_1)$ в представлении  


(12) является функцией одной переменной $\theta_1$, а уравнение (13) 


переходит в  


уравнение Лежандра ([9], с.\,125) и гармоническая функция 


представляется так 


\begin{equation} 


F(\theta_1)=C_0\rho^{m_0}_1P_{m_0}(\cos\theta_1)+ 


O(\rho^{m_0+1}_1), 


\end{equation} 


а ее нормальная производная при $m_0\in(0,1)$ принимает сингулярность в  


конической точке 


\begin{equation} 


\partial F/\partial n=\rho^{m_0-1}_1C_0d P_{m_0} 


(\cos\theta_1)/d\theta_1+O(\rho^{m_0}_1), 


\end{equation} 


где $P_{m_0}(\cos\theta_1)$ --- функции Лежандра, $C_0$ --- постоянная. 


 


Отметим, что в представления (18), (19) вошли только функции Лежандра  


первого рода как ограниченные при $\theta_1\in[-\pi,\pi)$. 


 


Подставляя выражения для гармонических функций и их производных типа  


(16), (17) в условия (3), записанные в локальных координатах $\rho_1$, 


$\theta_1$, $s_1$, 


полагая $\theta_1=\pi/2$ и приравнивая выражения при одинаковых 


степенях $\rho_1$ и гармониках  


нулю, получим однородную систему линейных алгебраических уравнений  


относительно произвольных постоянных. Приравнивая определитель такой  


системы нулю, получаем характеристическое уравнение для определения  


величины $m_0$ 


\begin{equation} 


\Delta_0(m_0)=0. 


\end{equation} 


Аналогично поступая с представлениями (18), (19), находим 


характеристическое  


уравнение 


\begin{equation} 


P_{m_0}(\cos\theta_1)=0. 


\end{equation} 


Таким образом, доказана  


 


\begin{lem} 


Если характеристические уравнения $(20)$, $(21)$ имеют корни 


$m_0\in(0,1)$, то гармонические функции и их нормальные производные 


имеют  


асимптотику $(16)$--$(19)$. 


\end{lem} 


 


Отметим, что множество корней $m_0\in(0,1)$ характеристических 


уравнений  


(20), (21) непустое ([10], с.\,320).  


\medskip 


 


{\bf 3.} {\it Разрешимость задачи}. Перейдем к общему случаю задачи 


сопряжения (1)--(3). 


 


\begin{thmnonum} 


Если в каждой точке особой линии выполняется условие 


$m=\pi/(2\pi-\varpi)<1$, 


а в каждой конической точке существуют корни $m_0\in(0,1)$ 


характеристических  


уравнений $(20)$, $(21)$, то задача сопряжения $(1)$--$(3)$ безусловно 


разрешима.  


\end{thmnonum} 


 


\begin{pf} 


Рассмотрим сначала случай одной конечной области $V_1$ с 


граничной поверхностью $S=\bigcup\limits^{n_1}_{j=1}S_j$, содержащей 


конические точки $O_k$, $k=\ov{1,m_1}$, 


и гладкие замкнутые непересекающиеся особые линии 


$L_j=S_j\cap S_{j+1}$, $j=\ov{1,n_1-1}$, $O_k\notin L_j$. Поверхности 


$S_j$ определяются уравнениями  


$f_j(x,y,z)=0$, а поверхности, натянутые на контуры $L_j$,


описываются уравнениями $f_{0j}(x,y,z)=0$. 


 


Представим гармонические функции в виде 


\begin{gather} 


F_0(M)=F_{01}(M)+F_{02}(M),\quad M\in V_0, \\ 


F_1(M)=F_{11}(M)+F_{12}(M),\quad M\in V_1, 


\end{gather} 


где первые слагаемые реализуют асимптотику (6), (16)--(19) в 


зависимости от  


типа негладкости, а вторые являются ограниченными непрерывными  


функциями в своих областях определения. 


 


Множество точек гладкости поверхности $S$ является кусочно-гладкой  


поверхностью. Представим первые слагаемые посредством потенциалов 


простого и двойного слоя, которые являются гармоническими исчезающими 


на 


бесконечности функциями ([11], с.\,398), 


\begin{multline} 


F_{i1}(M)=\int\limits_S\bigg(\prod^{n_1-1}_{j=1} 


\mu^{(1)}_{i j 1}(N)\bigg)\bigg(\prod^{p_1}_{k=1} 


\mu^{(2)}_{i k 1}(N)\bigg)R^{-1}d s+ \\ +\iint\limits_S\bigg( 


\prod^{n_1-1}_{j=1}\mu^{(3)}_{i j 1}(N)\bigg)\bigg( 


\prod^{p_1}_{k=1}\mu^{(4)}_{i k1}(N)\bigg)\partial R^{-1}/ 


\partial n\,d s, 


\end{multline} 


где $N=N(\xi,\eta,\chi)\in S$, 


$R=((x-\xi)^2+(y-\eta)^2+(z-\chi)^2)^{1/2}$, $i=0,1$, 


производные берутся по направлению внешней нормали к $S$. 


 


Для особых линий берем плотности в виде 


\begin{align} 


\mu^{(3)}_{i j1}(N)&=(-1)^{i+1}/(2\pi)R^{m_j}_j(N) 


(C_{i1j}\cos(m_j T_j(N))+C_{i2j}\sin(m_j T_j(N))),\\ 


\mu^{(1)}_{i j1}(N)&=(-1)^{i+1}m_j/(2\pi)R^{m_j-1}_j(N) 


(C_{i2j}\cos(m_j T_j(N))- 


C_{i1j}\sin(m_j T_j(N))), 


\end{align} 


где 


\begin{gather*} 


R_j(N)=((f_{0j}(N)/|\Grad f_{0j}(N)|)^2+ 


Q^2_j(N))^{1/2},\\ 


T_j(N)=\arcsin(f_{0j}(N)/(|\Grad f_{0j}(N)|R_j(N))),\\ 


Q_j(N)=f_j(N)/|\Grad f_j(N)|+(1-(\bold G_{2j}(N), 


\bold G_{0j}(N))^2f_{0j}(N)/|\Grad f_{0j}(N)|) 


(\bold G_{0j}(N),\bold G_{2j}(N)), \\ 


\bold G_j(N)=\Grad f_j(N)/|\Grad f_j(N)|, \\ 


\bold G_{0j}(N)=\Grad f_{0i}(N)/|\Grad f_{0j}(N)|, \\ 


\bold G_{2j}(N)=\big[[\bold G_j(N),\bold G_{j+1}(N)]/ 


|[\bold G_j(N),\bold G_{j+1}(N)]|,\bold G_{0j}(N)\big]. 


\end{gather*} 


Для конических точек имеем 


\begin{align} 


\mu^{(4)}_{i k1}(N)&=(-1)^{i+1}/(2\pi)R^{m_k}_k(N) 


\sum^K_{n=0}(D_{i k n1}(T_k(N))\cos(n\pi S_k(N)/s_{k0})+\notag \\ 


&\quad +D_{i k n2}(T_k(N))\sin(n\pi S_k(N)/s_{k0})), \\ 


% 


\mu^{(2)}_{i k1}(N)&=(-1)^{i+1}/(2\pi)R^{m_k-1}_k(N) 


\sum^K_{n=0}(d D_{i k n1}(T_k(N))/d T_k(N) 


\cos(n\pi S_k(N)/s_{k0})+\notag\\ 


&\quad +d D_{i k n2}(T_k(N))/d T_k(N)\sin(n\pi S_k(N)/s_{k0})), 


\end{align} 


где 


\begin{gather*} 


R_k(N)=((\xi-x_k)^2+(\eta-y_k)^2+(\chi-z_k)^2)^{1/2},\\ 


T_k(N)=\arcsin(l_{0k}(N)(\xi-x_k)+m_{0k}(N)(\eta-y_k)+n_{0k}(N) 


(\chi-z_k)), \\ 


S_k(N)=\int^{\varphi_k}_0((d l_{0k}(N)/d\varphi_k)^2+ 


(d m_{0k}(N)/d\varphi_k)^2+(d n_{0k}(N)/d\varphi_k)^2)^{1/2}d\varphi_k, 


\end{gather*} 


$\varphi_k=\varphi_k(\xi,\eta)$, $O_k=O_k(x_k,y_k,x_k)$, $l_{0k}(N)$, 


$m_{0k}(N)$, $n_{0k}(N)$ --- координаты орта образующей $\bold 


r_{0k}(N)$, $D_{i k n q}$ --- известные 


функции в представлении (15). 


 


Аналогичные формулы в частном случае асимптотики (18),(19) 


\begin{align*} 


\mu^{(4)}_{i k1}(N)&=(-1)^{i+1}/(2\pi)R^{m_k}_k(N)C_{0k} 


P_{m_k}(\cos T_{1k}(N)),\\ 


\mu^{(2)}_{i k1}(N)&=(-1)^{i+1}/(2\pi)R^{m_k-1}_k 


(N) C_{0k}d P_{m_k}(\cos T_{1k}(N))/d T_{1k}(N), 


\end{align*} 


где $T_{1k}(N)=T_k(N)+\beta$, $P_{m_k}$ --- функции Лежандра. 


 


Подставляя гармонические функции (22), (23) с учетом (24)--(28)  


в граничные условия (3), убеждаемся, что они удовлетворяются, а функции  


и их производные имеют асимптотику (6), (16)--(19). При этом учитываем, 


что  


согласно леммам 1, 2\; $m_j,m_k\in(0,1)$ и граничные значения 


потенциалов простого и 


двойного слоя в точках гладкости поверхности сопряжения определяются  


известными значениями [11], и переходим в них к локальным координатам  


согласно (7), (11). 


 


Вторые слагаемые гармонических функций (22), (23) записываем  


потенциалами простого и двойного слоя с неизвестными плотностями 


$\mu_{i2q}(N)$, $q=1,2$, 


\begin{gather} 


F_{02}(M)=F_{00}(M)+\iint\limits_S \mu_{021}(N)R^{-1}d s+ 


\iint\limits_S \mu_{022}(N)\partial R^{-1}/ 


\partial n\, d s,\\ 


F_{12}(M)=F^*_{11}(M)+\iint\limits_S \mu_{121}(N)R^{-1}d s+ 


\iint\limits_S \mu_{122}(N)\partial R^{-1}/ 


\partial n\,d s, 


\end{gather} 


где $F_{00}(M)$, $F^*_{11}(M)$ --- заданные гармонические слагаемые,  


обусловленные физической сущностью задачи [3]. 


 


Для нахождения плотностей подставляем (22), (23), (29), (30) в  


граничные условия (1), (2) и полагаем ввиду произвольности плотностей 


$$ 


\mu_{121}(N)=\mu_{021}(N),\quad 


\mu_{122}(N)=\mu_{022}(N). 


$$ 


В результате находим 


$$ 


\mu_{022}(N)=h_0(N), 


$$ 


а для плотности $\mu_{021}(N)$ получаем интегральное уравнение 


\begin{equation} 


\mu_{021}(N_0)-(\sigma/(2\pi))\partial\bigg(\iint\limits_S 


\mu_{021}(N)R^{-1}d s\bigg)/\partial n=g_0(N_0), 


\end{equation} 


где $N_0\in S$, 


\begin{gather*} 


h_0(N)=(F^-_{01}(N)+F^+_{00}(N)-(F^+_{11}(N)+F^{*+}_{11}(N)))/(4\pi),\\ 


% 


g_0(N_0)=(\gamma(\partial F^-_{01}/\partial n+\partial F^-_{00} 


/\partial n)-(\partial F^+_{11}/\partial n+\partial F^{*+}_{11} 


/\partial n))/(2\pi(1+\gamma))- \\ -(\sigma/(2\pi))\partial\bigg( 


\iint\limits_S h_0(N)\partial R^{-1}/\partial n\,d s 


\bigg)/\partial n,\\ 


% 


F_{i1}=F_{i1}(N_0),\ \ F_{00}=F_{00}(N_0),\ \  


F^*_{11}=F^*_{11}(N_0),\ \ \sigma=(1-\gamma)/(1+\gamma), 


\ \ \gamma=\lambda_0/\lambda_1. 


\end{gather*} 


Правая часть уравнения (31) ввиду выполнения условий (3) является  


непрерывной ограниченной функцией на $S$, поэтому это уравнение с 


полярным  


ядром ([11], с.\,279). 


 


Однородное союзное уравнение 


$$ 


\mu^*_{021}(N_0)-(\sigma/(2\pi))\iint\limits_S 


\mu^*_{021}(N)\partial R^{-1}/\partial n\,d s=0 


$$ 


можно трактовать как выполнение условия (1) для функции 


$$ 


F^*(M)=\Gamma\iint\limits_S \mu^*_{021}(N) 


\partial R^{-1}/\partial n\,d s, 


$$ 


причем $\Gamma=\gamma$, если $M\in V_0$, и $\Gamma=1$, если $M\in V_1$. 


Она является ввиду (1) непрерывной гармонической функцией во всем 


пространстве, принимающей постоянное значение ввиду (2) на $S$ 


и обращающейся в нуль на бесконечности. Следовательно, эта функция 


тождественно равна нулю ([11], с.\,374). 


 


В случае нескольких непересекающихся конечных односвязных областей  


$V_i$, $i=\ov{1,k_0}$, гармонические функции $F_i(M)$, $M\in V_i$, 


строим аналогично (23), (24)  


для каждой области. Гармоническую функцию $F_0(M)$, $M\in V_0$, берем в 


виде (24),  


где 


$$ 


F_{01}(M)=\sum^{k_0}_{i=1}F_{0i1}(M),\quad 


F_{02}(M)=\sum^{k_0}_{i=1}F_{0i2}(M); 


$$ 


здесь слагаемые $F_{0i1}(M)$ строятся с учетом негладкостей 


соответствующих  


поверхностей $S_i$ аналогично (24)--(28), а слагаемые $F_{0i2}(M)$ 


представляются в 


виде (29). 


\end{pf} 


 


Для конкретных физических задач рассмотренная задача дополняется 


условиями 


аналогично двумерному случаю [12], [13].  


 


\begin{thebibliography}{99} 


 


\bibitem{1} 


Денисюк И.Т. {\em Решение одной задачи сопряжения для составной области 


с угловыми точками на линиях раздела} // Изв. вузов. Математика. -- 1996. 


-- \No\,6. -- С.\,17--24.  


\bibitem{2} 


Денисюк И.Т. {\em Одна задача сопряжения аналитических функций в 


аффинно 


преобразованных областях с кусочно-гладкими границами} // Изв. вузов. 


Математика. -- 2000. -- \No\,6. -- С.\,70--74. 


\bibitem{3} 


Седов Л.И. {\em Механика сплошной среды}. Т.\,1. -- М.: Наука, 1973. -- 


536  с. 


\bibitem{4} 


Седов Л.И. {\em Механика сплошной среды}. Т.\,2. -- М.: Наука, 1973. -- 


568  с.  


\bibitem{5} 


Денисюк И.Т. {\em Напряженное состояние вблизи особой линии поверхности 


раздела сред} // Изв. РАН. Механ. тверд. тела. -- 1995. -- 


\No\,5. -- С.\,64--70.  


\bibitem6{} 


Норден А.П. {\em Краткий курс дифференциальной геометрии}. -- М.: 


ГИФМЛ, 1958, 


-- 244 с. 


\bibitem{7} 


Фихтенгольц Г.М. {\em Курс дифференциального и интегрального 


исчисления}. 


Учеб. пособие. Т.\,3. -- М.: Наука, 1969. -- 656 с.  


\bibitem{8} 


Петровский И.Г. {\em Лекции по теории обыкновенных дифференциальных 


уравнений}. -- М.: Наука, 1979. -- 280 с.  


\bibitem{9} 


Бейтмен Г., Эрдейи Г. {\em Высшие трансцендентные функции}. Ч.\,1. -- 


М.:  


Наука, 1974. -- 296 с. 


\bibitem{10} 


Партон В.З., Перлин П.И. {\em Методы математической теории упругости}. 


Учеб. пособие. -- М.: Наука, 1981. -- 688 с. 


\bibitem{11} 


Владимиров В.С. {\em Уравнения математической физики}. 


-- М.: Наука, 1971. -- 512 с.  


\bibitem{12} 


Денисюк И.Т. {\em Термоупругость изотропной пластинки с угловыми 


включениями} // Изв. РАН. Механ. тверд. тела. -- 1999. -- 


\No\,2. 


-- C.\,148--155.  


\bibitem{13}  


Денисюк И.Т. {\em Одна модель тонких упругих включений в изотропной  


пластинке} // Изв. РАН. Механ. тверд. тела. -- 2000. -- 


\No\,4. -- С.\,140--148. 


 


\end{thebibliography} 


 


\vskip 2cm 


 


\post{\it Луцкий государственный }{\it Поступила} 


\post{\it технический университет $($Украина$)$}{28.11.2000} 


 


\label{end} 


\end{document} 


